\section{A quotient remembers which identifications were made}
\label{r:chapter}

We now have a language in which a quotient is an operation on a relation,
and a comparison of quotients is an operation on a correspondence.  We will
use that language once, in a pointed Barr-exact Mal'cev category
$\mathcal C$.  Groups and modules will enter afterwards by verifying these
axioms.  Neither commutative addition of arrows nor an invariant called
dimension is used in this chapter.

There are three stages.  First we record a normal series without forgetting
its quotient maps.  Next we compare two such records by a rectangular array
of correspondences.  Finally simplicity prevents the comparison from
breaking a factor into smaller nonzero factors.  The last step is
Jordan--H\"older; the rectangle, rather than a numerical count, does the work.

\subsection{The data in a normal series}

\begin{definition}[Normal series with its quotient maps]
\label{r:normal-series}
A finite normal series of $X$ consists of a chain of subobjects and chosen
regular epimorphisms
\[
  0=X_0\lhook\joinrel\longrightarrow X_1
  \lhook\joinrel\longrightarrow\cdots
  \lhook\joinrel\longrightarrow X_n=X,
  \qquad q_i:X_i\twoheadrightarrow S_i,
\]
such that the inclusion $X_{i-1}\hookrightarrow X_i$ is a kernel of $q_i$.
The objects $S_i$ are its \emph{factors}.  Repetitions of subobjects and zero
factors are permitted.  Two series are \emph{equivalent} when their lists of
factors are isomorphic after a permutation.  We also use equivalence
\emph{after deleting zero factors}, stating explicitly when we do so.
\end{definition}

The kernel square in this definition is
\[
\begin{tikzcd}[column sep=large,row sep=large]
X_{i-1}\arrow[r,hook]\arrow[d] & X_i\arrow[d,two heads,"q_i"]\\
0\arrow[r] & S_i .
\end{tikzcd}
\tag{\ref{r:normal-series}a}
\]
It is a pullback.  In particular, the normality assertion is relative to
$X_i$, not to all of $X$.  For a group this allows a subgroup normal in its
successor without demanding that it be normal in the ambient group.

We keep $q_i$ because the kernel alone need not determine a quotient in
this generality.  A familiar implication from groups, such as ``a regular
epimorphism with zero kernel is an isomorphism'', is not inserted as an
unstated axiom.  All our comparisons retain the kernel pair of the actual
map, which is the full record of its identifications.

Here and below, composition of relations is function-like: $SR$ means
first $R$, then $S$.  For a relation $r:X\rightsquigarrow X$ and a subobject
$U\hookrightarrow X$, write $rU$ for the image of the restriction of $r$
whose input lies in $U$.  Thus $rX$ is its output support and $r0$ its zero
class.  A symmetric transitive relation is a \emph{partial equivalence
relation}; it is an equivalence relation on its support, by
Lemma~\ref{d:partial-equivalence}.  Its support quotient will be written
\[
                    Q(r)=(rX)/r.
\]
The slash is shorthand for the effective quotient of the restricted
equivalence relation, not a new operation on objects without relations.

\begin{lemma}[A series as a sequence of identifications]
\label{r:series-relations}
For $n\geq1$, the data of Definition~\ref{r:normal-series}, up to the unique isomorphisms
allowed for quotient objects, are equivalently a list of partial equivalence
relations $r_1,\ldots,r_n$ on $X$ satisfying
\begin{equation}
\label{r:boundary}
 r_1 0=0,\qquad r_iX=r_{i+1}0\ (1\leq i<n),\qquad r_nX=X.
\end{equation}
The associated factors are $Q(r_i)$. The empty series of $0$ is treated
separately; it has no endpoint relations to name.
\end{lemma}

\begin{proof}
Let $a_i:X_i\hookrightarrow X$ be the composite inclusion.  Regarding
arrows as their graphs, put
\[
       r_i=a_iq_i^{\mathrm{op}}q_i a_i^{\mathrm{op}}.
\]
On $X_i$, this is the kernel pair of $q_i$; it relates precisely the two
inputs that $q_i$ identifies.  Extending it along the monomorphism $a_i$
makes it a partial equivalence relation on $X$.  Its support is $X_i$.
The zero arrow into $X_i$ is sent to the zero arrow into $S_i$, so taking
the zero class of that kernel pair gives precisely the pullback defining
$\ker q_i$.  Therefore
\[
                 r_iX=X_i,\qquad r_i0=X_{i-1}.
\]
These equalities give~\eqref{r:boundary}.  Since $q_i$ is a regular
epimorphism, it is the effective quotient of its kernel pair, and
$Q(r_i)\cong S_i$ canonically.

Conversely, set $X_i=r_iX$.  The zero class $r_i0$ is a subobject of $r_iX$.
Thus the identities~\eqref{r:boundary} make the $X_i$ a chain, beginning
with $X_0=0$ and ending with $X_n=X$.  By exactness the restricted relation
$r_i$ has an effective quotient $q_i:X_i\twoheadrightarrow Q(r_i)$.
The same kernel-pair calculation says that $\ker q_i=r_i0=X_{i-1}$.
These two constructions are inverse because an effective quotient is
determined, up to a unique compatible isomorphism, by its kernel pair.
\end{proof}

This change of language has retained more than the chain.  It retains
which pairs of points become equal at each step.  That is the information
needed to compare two quotient constructions in a category where the zero
class alone may not determine the entire relation.

\subsection{What a genuine refinement means}

To refine a series, we refine one of its factors and pull that refinement
back.  Suppose $q:D\twoheadrightarrow Q$ is one quotient step, with kernel
$K\hookrightarrow D$, and suppose
\[
 0=Q_0\hookrightarrow Q_1\hookrightarrow\cdots\hookrightarrow Q_t=Q,
 \qquad e_j:Q_j\twoheadrightarrow Z_j,\quad\ker e_j=Q_{j-1},
\]
is a normal series of its factor.  Form $P_j=D\times_Q Q_j$.  We obtain
\[
\begin{tikzcd}[column sep=large,row sep=large]
K=P_0\arrow[r,hook]\arrow[d,two heads]
 &P_1\arrow[r,hook]\arrow[d,two heads,"p_1"]
 &\cdots\arrow[r,hook]&P_t=D\arrow[d,two heads,"q"]\\
0=Q_0\arrow[r,hook]&Q_1\arrow[r,hook]
 &\cdots\arrow[r,hook]&Q_t=Q .
\end{tikzcd}
\]
Each displayed comparison square is a pullback.  Pullback stability of
regular epimorphisms makes $p_j:P_j\twoheadrightarrow Q_j$ a regular
epimorphism.  The refined quotient map at the $j$th step is
\begin{equation}
\label{r:refined-map}
                  e_jp_j:P_j\twoheadrightarrow Z_j.
\end{equation}
Its kernel is $P_{j-1}$: its pullback over $0$ is
$P_j\times_{Q_j}Q_{j-1}=D\times_Q Q_{j-1}$.  Regular epimorphisms are
closed under composition, so this really is a normal series between $K$
and $D$.  Replacing any or all of the original steps in this way is what
we mean by a \emph{refinement}.
For the empty series of the zero object we adopt the harmless convention
that inserting a finite list of zero steps is a refinement.  This
convention concerns only how empty lists are represented; it introduces
no nonzero factor.

It would be insufficient to insert arbitrary kernel subobjects between
$K$ and $D$ while forgetting their relation to $q$.  We need the new
quotients to have arisen from the original quotient.  The following
saturation condition expresses exactly that requirement.

\begin{proposition}[Saturation is descent of a refinement]
\label{r:refinement-descent}
Let $r$ be a partial equivalence relation on $X$, let $D=rX$, and let
$q:D\twoheadrightarrow Q(r)$ be its quotient.  For $k\geq1$, lists of
partial equivalence relations $t_1,\ldots,t_k$ on $X$ satisfying
\begin{align}
 rt_jr&=t_j \quad(1\leq j\leq k), \label{r:saturation}\\
 t_1 0&=r0,\qquad t_jX=t_{j+1}0\ (1\leq j<k),\qquad
 t_kX=rX \label{r:refinement-boundary}
\end{align}
correspond to normal series of $Q(r)$ with exactly $k$ steps. Under this correspondence, their
support quotients $Q(t_j)$ are the factors of the pulled-back refinement
of $q$.
\end{proposition}

\begin{proof}
The equality $rt_jr=t_j$ first implies that the support of $t_j$ lies
inside $D$, so we can work on $D$.  Write $E=q^{\mathrm{op}}q$ for the
kernel pair on $D$; it is the restriction of $r$.  The saturated descent
lemma, Proposition~\ref{d:saturated-descent}, gives mutually inverse
operations
\[
       T_j=q\,t_j\,q^{\mathrm{op}},\qquad
       t_j=q^{\mathrm{op}}T_jq.
\]
The first is the direct image relation on $Q(r)$; the second is its
pullback along $q\times q$.  Their being inverse uses $qq^{\mathrm{op}}
=1_{Q(r)}$ and $Et_jE=t_j$.  That lemma also proves that symmetry and
transitivity descend, so each $T_j$ is a partial equivalence relation.

We check the endpoints, since those checks are what make this a descent
of \emph{series}.  Saturation and reflexivity of $E$ on $D$ give
$t_jE=t_j=Et_j$.  The inverse image of the zero subobject of $Q(r)$ is
$E0$.  Therefore, as subobjects obtained by regular images,
\[
             T_j0=q\,t_jE0=q\,t_j0,
       \qquad T_jQ(r)=q(t_jD).
\]
Now $q(E0)=0$ and $qD=Q(r)$.  Applying these formulas to
\eqref{r:refinement-boundary} gives
\[
 T_10=0,\qquad T_jQ(r)=T_{j+1}0,\qquad T_kQ(r)=Q(r).
\]
Lemma~\ref{r:series-relations} therefore supplies a normal series of
$Q(r)$.

In the other direction, begin with a normal series of $Q(r)$ and its
relations $T_j$.  Pulling back the relations gives $t_j=q^{\mathrm{op}}T_jq$,
which are saturated and have precisely the support and zero-class
boundaries in~\eqref{r:refinement-boundary}.  Concretely, if $Q_j$ is the
support of $T_j$, the support of $t_j$ is $P_j=D\times_{Q(r)}Q_j$.
If $e_j:Q_j\twoheadrightarrow Z_j$ is its factor quotient, the kernel pair
of $e_jp_j$ on $P_j$ is $t_j$.  As $e_jp_j$ is a regular epimorphism,
its codomain $Z_j$ is the effective quotient of $t_j$.  This gives the
unique compatible isomorphism $Q(t_j)\cong Z_j$ and identifies the
quotient maps as well as the objects.
\end{proof}

The proposition is a small descent theorem.  The equality $rt_jr=t_j$
says that changing either input by an identification already imposed by
$r$ does not change the new relation.  Exactly then the relation can be
defined downstairs, independently of the choice of a representative.
The proof above replaces that verbal test by pullbacks, images, and
unique quotient maps.

\subsection{The comparison rectangle}

Let $r,s$ be partial equivalence relations on the same object.  Their
composite $rs$ need not be an equivalence relation on all of $X$.  It is
nevertheless a relation, so the Mal'cev theorem applies to it.  This
produces the identities that will make our rectangle fit together.

\begin{lemma}[Two endpoint identities]
\label{r:endpoint-identities}
For partial equivalence relations $r,s$ on $X$,
\begin{equation}
\label{r:rs-idempotent}
                  rsrs=rs,
\end{equation}
and
\begin{equation}
\label{r:endpoint-equalities}
                  rsr0=rs0,\qquad rsrX=rsX.
\end{equation}
\end{lemma}

\begin{proof}
The relation $P=rs$ is difunctional by Proposition~\ref{d:difunctional}, so
$PP^{\mathrm{op}}P=P$.  Symmetry and idempotence of $r,s$ reduce the left
side to
\[
           rssrrs=rsrs,
\]
which proves~\eqref{r:rs-idempotent}.

Every subobject in a pointed category contains its zero arrow.  In
particular $0\subseteq r0$ and $0\subseteq s0$.  Monotonicity of relational
image and~\eqref{r:rs-idempotent} now give
\[
 rs0\subseteq rsr0\subseteq rsrs0=rs0.
\]
Here the second inclusion comes from $r0\subseteq rs0$, itself obtained
by applying $r$ to $0\subseteq s0$.  For supports, $sX\subseteq X$ gives
\[
        rsX=rsrsX\subseteq rsrX\subseteq rsX;
\]
the final inclusion uses $rX\subseteq X$.  Both chains give equality.
\end{proof}

Now take two normal series of $X$ of positive lengths $m,n$, encoded by
$r_1,\ldots,r_m$ and $s_1,\ldots,s_n$.  At the crossing of the $i$th row
and $j$th column, define
\begin{equation}
\label{r:cells}
       r_{ij}=r_is_jr_i,\qquad
       s_{ji}=s_jr_is_j,\qquad P_{ij}=r_is_j.
\end{equation}
The two cell relations are the two equivalence relations associated with
the correspondence $P_{ij}$:
\[
       P_{ij}P_{ij}^{\mathrm{op}}=r_{ij},\qquad
       P_{ij}^{\mathrm{op}}P_{ij}=s_{ji}.
\]
Our composition convention matters here: $s_{ji}$ is on the input support
of $P_{ij}$, and $r_{ij}$ on its output support.

\begin{proposition}[Every row and column is a genuine refinement]
\label{r:rectangle-refines}
For fixed $i$, the list $r_{i1},\ldots,r_{in}$ satisfies
Proposition~\ref{r:refinement-descent} for $r_i$.  For fixed $j$, the list
$s_{j1},\ldots,s_{jm}$ satisfies that proposition for $s_j$.
\end{proposition}

\begin{proof}
We verify every condition for a row.  Symmetry follows by reversing the
word $r_is_jr_i$.  Idempotence follows from
Lemma~\ref{r:endpoint-identities}:
\[
 r_{ij}r_{ij}=r_is_jr_is_jr_i=r_is_jr_i=r_{ij}.
\]
Thus each cell is symmetric and transitive.  Its saturation is
\[
       r_i r_{ij}r_i=r_i^2s_jr_i^2=r_{ij}.
\]
For the first zero class, use $s_10=0$:
\[
            r_{i1}0=r_is_1r_i0=r_is_10=r_i0.
\]
For an intermediate join, use $s_j0=s_{j-1}X$:
\[
\begin{split}
 r_{ij}0
 &=r_is_jr_i0=r_is_j0
  =r_is_{j-1}X\\
 &=r_is_{j-1}r_iX=r_{i,j-1}X
                  \qquad(2\leq j\leq n).
\end{split}
\]
For the last support, use $s_nX=X$:
\[
                 r_{in}X=r_is_nr_iX=r_is_nX=r_iX.
\]
Every equality that removes the trailing $r_i$ is one of
\eqref{r:endpoint-equalities}.  We have proved saturation and all boundary
conditions, so Proposition~\ref{r:refinement-descent} supplies a
pulled-back refinement of the original quotient.  Interchanging $r$ and
$s$ proves the column assertion.
\end{proof}

The boundary identities are the gluing conditions.  They say that the
support of one cell is precisely the zero class required to start the
next cell.  Without these identities we would have an assortment of
quotients, not a refinement of a filtration.

\begin{proposition}[The two readings of a cell have the same quotient]
\label{r:cell-isomorphism}
For every $(i,j)$ there is a canonical isomorphism
\begin{equation}
\label{r:cell-iso}
                       Q(s_{ji})\xrightarrow{\ \sim\ }Q(r_{ij}).
\end{equation}
It is induced by the correspondence $P_{ij}$ and characterized by its
two projections, not chosen by comparing lengths or other invariants.
\end{proposition}

\begin{proof}
Let $C_{ij}$ be the object representing $P_{ij}\hookrightarrow X\times X$.
Factor its two projections through their image supports:
\[
 C_{ij}\xrightarrow{a_{ij}}L_{ij}\hookrightarrow X,
 \qquad C_{ij}\xrightarrow{b_{ij}}R_{ij}\hookrightarrow X.
\]
Both $a_{ij}$ and $b_{ij}$ are regular epimorphisms.  The input relation
$P_{ij}^{\mathrm{op}}P_{ij}$ is $s_{ji}$, so $L_{ij}=s_{ji}X$; similarly
$R_{ij}=r_{ij}X$.  Write their effective quotient maps as
$\lambda_{ij}:L_{ij}\twoheadrightarrow Q(s_{ji})$ and
$\rho_{ij}:R_{ij}\twoheadrightarrow Q(r_{ij})$.

The support-quotient comparison theorem, Theorem~\ref{d:comparison},
applies to $P_{ij}$ and gives the unique isomorphism $\theta_{ij}$ for
which
\[
\begin{tikzcd}[column sep=large,row sep=large]
L_{ij}\arrow[d,two heads,"\lambda_{ij}"']
 &C_{ij}\arrow[l,two heads,"a_{ij}"']
          \arrow[r,two heads,"b_{ij}"]
 &R_{ij}\arrow[d,two heads,"\rho_{ij}"]\\
Q(s_{ji})\arrow[rr,"\theta_{ij}"',"\sim"]
 &&Q(r_{ij})
\end{tikzcd}
\]
commutes, meaning
$\theta_{ij}\lambda_{ij}a_{ij}=\rho_{ij}b_{ij}$.
The cited theorem proves existence and invertibility by effective
quotients and difunctionality.  Its uniqueness also follows directly
because $\lambda_{ij}a_{ij}$ is an epimorphism.  The displayed computation
of the two kernel relations identifies its conclusion with
\eqref{r:cell-iso}.
\end{proof}

This is the categorical butterfly.  It is the same correspondence read
from its two ends, with each end quotiented by precisely the ambiguity
introduced by that correspondence.  No picture of four wings is needed
to remember the construction: form $P$, compute $P^{\mathrm{op}}P$ and
$PP^{\mathrm{op}}$, and take their support quotients.

\begin{theorem}[Schreier refinement in an exact Mal'cev category]
\label{r:schreier}
\label{j:schreier}
In a pointed Barr-exact Mal'cev category, any two finite normal series of
an object have equivalent refinements in the pullback sense specified
above.  The refinements may contain repeated subobjects and zero factors.
\end{theorem}

\begin{proof}
Suppose first that both series have positive numbers $m,n$ of steps.
For the first series, concatenate the $m$ rows supplied by
Proposition~\ref{r:rectangle-refines}.  Their endpoints are the original
successive subobjects, so they fit together into a normal series of $X$.
For the second series, concatenate the $n$ columns.  Both refined lists
have $mn$ factors.  Their indexing orders differ, but the transposition
$(i,j)\leftrightarrow(j,i)$ pairs them.  Proposition~\ref{r:cell-isomorphism}
gives the isomorphism for each paired cell.

An empty series can occur only for $X=0$.  Every subobject of $0$ is
isomorphic to $0$: a monomorphism $A\to0$ has its unique map $0\to A$ as
a right inverse, and monicity makes it a left inverse as well.  A regular
epimorphism out of $0$ also has zero codomain, since it is at once monic
and a regular epimorphism.  Thus every factor of every series of $0$ is
zero.  We may insert zero steps to give both series the same positive
number of zero factors, or treat two empty lists as already equivalent.
These are refinements by the zero object's zero series.  This resolves
the degenerate case without assuming a nonempty comparison rectangle.
\end{proof}

The theorem does not produce a single nested chain containing the terms
of both original chains.  Two different lines in a two-dimensional vector
space already prevent that: neither line contains the other.  It produces
one refinement of each chain and a matching between their quotient
factors.  This distinction matters even in the most elementary example.

\subsection{Simplicity removes the remaining freedom}

\begin{definition}[Quotient-simple object]
\label{j:simple}
A nonzero object $S$ is \emph{quotient-simple} if every regular
epimorphism $e:S\twoheadrightarrow T$ is either an isomorphism or has
zero codomain.  A composition series is a finite normal series all of
whose factors are quotient-simple.
\end{definition}

The condition concerns the arrow $e$.  The weaker phrase ``its codomain
is isomorphic to $S$ or to $0$'' would not exclude a proper self-quotient.
The definition above does exclude it.  In the group and module instances,
it agrees with the usual definition of a nontrivial simple object.

\begin{lemma}[A simple factor cannot be broken into smaller factors]
\label{j:simple-refinement}
Every finite normal series of a quotient-simple object $S$ has exactly
one nonzero factor, isomorphic to $S$.  Its other factors are zero.
\end{lemma}

\begin{proof}
The last quotient is a regular epimorphism $e:S\twoheadrightarrow T$.
If $T=0$, its kernel is all of $S$, so the preceding subobject is $S$.
Remove this last, zero-factor step and consider the shorter series.
Because the series is finite and begins at $0\neq S$, this removal
process must eventually reach a last step whose preceding subobject is
proper in $S$.

For that step the quotient cannot have zero codomain: a map to $0$ has
kernel $S$.  Simplicity therefore makes the quotient map an isomorphism.
Its kernel is $0$, so every preceding subobject of the chain is $0$.
Every regular quotient of $0$ is $0$, as checked in the proof of
Theorem~\ref{r:schreier}.  We have one quotient isomorphic to $S$, zero
quotients before it, and the zero quotients stripped from its end.
\end{proof}

\begin{theorem}[Jordan--H\"older]
\label{j:jordan-holder}
Let $\mathcal C$ be a pointed Barr-exact Mal'cev category.  If an object
$X$ has two finite composition series, then their factor lists agree,
including multiplicity, up to permutation and isomorphism.  In particular,
the two series have the same number of factors.
\end{theorem}

\begin{proof}
Apply Theorem~\ref{r:schreier} to obtain equivalent refinements.  Each
refinement is obtained from normal series of the original factors, pulled
back along the original quotient maps.  By Lemma~\ref{j:simple-refinement},
each such factor series contributes exactly its original simple factor
and some zero factors.  Proposition~\ref{r:refinement-descent} identifies
these with the factors of the pulled-back refinement.

Consequently, deleting zero factors from the first refined list recovers
the first original list, and deleting them from the second recovers the
second original list.  An isomorphism cannot identify a nonzero object
with $0$, so the matching between refined lists restricts to a matching
between their nonzero lists.  This is the desired permutation and set of
factor isomorphisms.  The equality of lengths is a consequence of that
bijection, not an input to its construction.
\end{proof}

No existence assertion is hidden in this theorem.  We assumed two finite
composition series and compared them.  A category may contain objects
with no finite composition series.  Even when maximal or minimal
subobjects can be selected, a termination argument is still required
before one can assert existence.  The finite sequences used here are
already supplied, so neither an ascending nor a descending chain
condition enters the uniqueness proof.

\subsection{Factor data becomes an invariant only after the comparison}

Now, and only now, a numerical or algebraic record of the factors becomes
legitimate.  Call an object \emph{composition-finite} when it has a finite
composition series, allowing $0$ with its empty series.  Write $\mu(X)$
for the finite multiset of isomorphism classes of its simple factors.
Theorem~\ref{j:jordan-holder} proves that $\mu(X)$ is independent of the
chosen series.  It does not reconstruct $X$: distinct extension data can
give the same multiset.

\begin{proposition}[Splicing an extension]
\label{j:extension-additivity}
\label{j:splicing}
Suppose $e:E\twoheadrightarrow Q$ is a regular epimorphism with kernel
$K\hookrightarrow E$.  If $K$ and $Q$ are composition-finite, then so is
$E$, and
\begin{equation}
\label{j:mu-additive}
                         \mu(E)=\mu(K)+\mu(Q).
\end{equation}
Here $+$ means disjoint union of finite multisets, including multiplicity.
\end{proposition}

\begin{proof}
Choose composition series for $K$ and $Q$.  Pull the series of $Q$ back
along $e$.  The construction~\eqref{r:refined-map} gives a normal series
from $K$ to $E$ whose factors are exactly those of $Q$.  Prepend the
chosen composition series from $0$ to $K$.  Every factor in the resulting
series is simple, so it is a composition series of $E$.  Its factor list
is the concatenation of the two chosen lists.  Jordan--H\"older makes
that calculation independent of all the choices and proves
\eqref{j:mu-additive}.
\end{proof}

This proposition establishes closure under these extensions.  It does
not assert, without an additional argument, that every subobject or every
regular quotient of a composition-finite object is composition-finite in
our weakest categorical setting.  The direction used in the proof is
precisely the direction stated.

For the next statement assume that the isomorphism classes of simple
objects under discussion form a set $\Sigma$, for example by working in
an essentially small category.  The free commutative monoid
$\mathbb N^{(\Sigma)}$ consists of finitely supported functions
$\Sigma\to\mathbb N$, with pointwise addition.  Thus $\mu(X)$ is an
element of this monoid.

\begin{corollary}[The universal additive record]
\label{j:universal-invariant}
\label{j:universal-monoid}
Let $M$ be a commutative monoid.  An arbitrary assignment $w:\Sigma\to M$
extends uniquely to an isomorphism invariant $v$ of composition-finite
objects satisfying
\[
 v(0)=0,\qquad v(E)=v(K)+v(Q)
\]
for every kernel--regular-quotient sequence $K\hookrightarrow
E\twoheadrightarrow Q$ whose terms are composition-finite.  The formula is
\[
                 v(X)=\sum_{[S]\in\Sigma}\mu(X)_{[S]}\,w([S]).
\]
\end{corollary}

\begin{proof}
Only finitely many summands are nonzero.  Jordan--H\"older makes the
formula well defined, and Proposition~\ref{j:extension-additivity} gives
the required additivity.  For uniqueness, every intermediate term of a
composition series is composition-finite, because its initial segment
is itself a composition series.  Repeatedly applying additivity to these
successive quotient maps forces $v(X)$ to equal the displayed sum.
\end{proof}

Taking every $w([S])=1$ defines length.  Taking values in a group gives
the corresponding universal group of additive invariants.  These are
consequences of factor comparison; using them to justify that comparison
would reverse the dependency and create circular reasoning.  The order
here is deliberate: quotient construction, refinement, factor matching,
then invariant.

\subsection{What is natural in the construction}

\begin{proposition}[Transport by a regular functor]
\label{r:functoriality}
Let $F:\mathcal C\to\mathcal D$ be a functor between pointed Barr-exact
Mal'cev categories that preserves finite limits, regular epimorphisms,
and a zero object.  Then $F$ transports the relation rectangle, its
refinements, and its cell isomorphisms to the corresponding constructions
in $\mathcal D$, up to their canonical structural isomorphisms.
\end{proposition}

\begin{proof}
Finite-limit preservation gives preservation of monomorphisms, products,
kernel pairs, kernels, and pullbacks.  If a map is factored as a regular
epimorphism followed by a monomorphism, its image under $F$ is again
such a factorization.  Uniqueness of image factorizations therefore shows
that $F$ preserves images.  Relational composition is an image of a
pullback projection, so $F(SR)$ identifies with $F(S)F(R)$.  Converse
relations are preserved by preservation of the product symmetry.

If $q$ is the effective quotient of an equivalence relation $E$, then
$F(q)$ is a regular epimorphism and its kernel pair is $F(E)$.  Hence it
is the effective quotient of $F(E)$.  These observations preserve every
operation in~\eqref{r:cells}, the support and zero-class equations, and
the pullback squares defining refinements.  Finally the image of the
commuting diagram in Proposition~\ref{r:cell-isomorphism} satisfies the
same characterization as the cell comparison in $\mathcal D$.  Its
uniqueness identifies the two isomorphisms.
\end{proof}

A regular functor need not preserve simple objects.  For instance,
forgetting a module's scalar action can reveal additive subgroups that
were not submodules.  The proposition transports the comparison
construction; a statement about composition series additionally requires
that the relevant simple factors remain nonzero and quotient-simple.
This distinction prevents a change of framework from silently forgetting
part of the theorem's information.

\subsection{Dependency and scope check}

The ordinary categorical proof is now complete.  The foundation needed
finite limits and regular images for relation composition, exactness for
the effective quotients, and the Mal'cev condition for difunctionality.
Pointedness supplied zero classes and kernel boundaries.  Saturated
descent made each row a refinement of the \emph{chosen} quotient map.
The support-quotient comparison matched each cell.  Simplicity then
removed the zero cells.  No invariant length, dimension, Grothendieck
group, or pre-existing Jordan--H\"older theorem entered the chain of
dependencies.

The comparison pattern is the classical categorical refinement argument
associated with Lambek; Knop's Appendix~B2 gives its exact Mal'cev
formulation and emphasizes the saturation requirement.  The proof here
has supplied the descent, boundary, and quotient-map checks explicitly;
it is not a claim of a new theorem.  The detailed primary reference is
\href{https://arxiv.org/pdf/math/0610552}{Friedrich Knop, \emph{Tensor
envelopes of regular categories}, Appendix B2}.

What has become general is a construction with stated hypotheses, not
the assertion that every category has composition factors.  The later
homotopical chapters will ask which quotient and refinement data survive
when equalities become coherent identifications.  The present theorem
is available there whenever its ordinary exact Mal'cev hypotheses are
actually verified; it does not, by terminology alone, assert a
Jordan--H\"older theorem for arbitrary higher objects.
